\chapter{The Machine Measures Its Own Measurement}

\section{Self-application}

The book has one climax left, and it is the second the computational gauge owns. The first, chapters ago, turned the
compiler into a meter and pointed it at its own act, reading the cost of measuring itself and staying bounded doing it.
That was the machine weighing its own working. This second climax goes one step further, to the thing weighing was always
leading toward: the machine does not merely meter its own act; it \emph{verifies} it. The self-hosting checker --- itself
a term the machine can hold and run --- is turned on the machine's own descriptions, and asked not what they cost but
whether they are \emph{the same}. And it answers.

Recall the shape of the first climax, because this one completes it. There, the meter read the overhead of the machine
accounting for its two descriptions of one object --- but it read the \emph{cost} of that accounting, its price in beats,
not the accounting itself. This chapter is the accounting itself. The checker that verifies the machine's terms is a term
of the same machine --- this is what self-hosting means, that the checker is written in the language it checks --- so it
can be turned on the machine's own claims about itself. Self-application is exactly this: the checker checking the
checker's own object. And the object it checks is the deepest claim the machine makes: that its two names for one thing
name one thing. Where the first climax billed the act, the second performs it.

The claim, stated plainly, is small and famous. The machine has two descriptions of a single object --- two ways of
writing what is, underneath, one thing. And it proves they are equal: not that they happen to agree, not that they are
close, but that they are \emph{the same}, by the principle that two propositions holding under exactly the same conditions
are the same proposition. The theorem is one line, and the machine checks it: the two descriptions denote one object; the
two names are one name. This is the machine measuring its own measurement in the strictest sense --- taking its own two
accounts of a thing and proving, by its own checker, that they were an account of \emph{one} thing all along. The
instrument that spent the book telling differences apart here proves an identity: that beneath two descriptions there was
never more than one object.

And now the fact that makes this the honest climax and not a grand one, and it is the fact the theory of self-reference
would least lead you to expect. A machine that checks its own claims risks the wall Turing~\cite{turing1936} and
G\"odel~\cite{godel1931} built: self-reference that diverges, a system asked about itself that cannot answer, a checker
that runs forever on its own object. This machine does not hit that wall, and the reason is that it asks the \emph{right}
question of itself --- not the undecidable ``do I halt?'' but the decidable ``are these two of my descriptions the
same?'' --- and asks it where the answer is settled. When the machine proves this deepest of its claims, and you ask what
the proof rested on, the answer is almost nothing: a single principle for identifying propositions that hold under the
same conditions, and beneath it the plainest ground there is, that a thing is what it is. No oracle, no completed
infinity, no leap the earlier chapters were careful never to take. The deepest identity the machine can prove is also the
cheapest --- it is bought with the emptiest ledger in the book, the same empty ledger the very first difference was
recorded on. Kleene's~\cite{kleene1952} recursion theorem guarantees that a self-hosting system \emph{has} such a fixed
point, a place where it can speak of itself without diverging; this machine's capstone is that fixed point, reached on
the plainest logic and owing almost nothing for it.

In the two verbs self-application is the machine's checker turned on its own working. The checker \emph{tanges} --- it is
the machine's own instrument of telling apart, the thing that decides whether two terms are distinct --- and here it is
aimed at the machine's own two descriptions. And what it finds, and proves, is a \emph{funge}: that the two descriptions
pool into one object, that beneath the two names there is a single thing. The self-application is a tange that proves a
funge --- the machine's tool for telling apart, used to prove that two of its own names never told apart anything real.
Telling-apart, run on the machine itself, certifies a sameness. What that self-verification \emph{costs}, and in what
currency, is the next section's.

\section{The witness bears the cost}

The last section proved the machine's deepest claim on nearly no assumptions. But proving is not free --- the first
climax established that; every proof the machine runs bills it heartbeats. So the proof of self-application, the witness
that the two descriptions are one, has a cost of its own. This section reads that cost, and finds in it the reading that
ties the whole book together.

Start with the cost itself, which is a subtraction. Proving that two descriptions denote one object costs the machine
something to elaborate --- and so does merely stating one of those descriptions, without proving anything. Take the cost
of the proof and subtract the cost of the mere description it proves, and what is left is the \emph{extra}: the price the
machine pays to \emph{prove} its two names are one, over and above the price of simply writing one of them down. That
difference is the overhead of self-application --- the surcharge on proving-over-stating, the cost the witness bears for
being a witness and not just a claim. The machine reads this overhead off its own meter, because by the first climax the
meter reads exactly such costs, and it reads it in the incompressible sense the third chapter gave: the overhead is the
part of the self-proof that cannot be shortened away, the irreducible price of the machine's self-reference.

And that is where the book's two climaxes meet. The first turned the compiler into a meter and pointed it at its own act;
the second turned the checker on its own truth and proved it. Here the meter of the first weighs the witness of the
second: the overhead of self-application, read in the machine's own heartbeats and scaled into its own units by its own
constant. The machine measures the cost of its own act of self-verification, in the currency it runs on. The two climaxes
were never two separate feats; they were a single act --- measure yourself, and verify yourself --- and this is the point
where they close on each other, the meter reading the price of the proof. The cost of the reading \emph{is} the reading:
the machine asked what it costs to know itself one thing, and the cost it read back is the coupling.

Now the reading that gathers the whole instrument into one. This overhead --- the price the machine pays to prove its own
account is a single thing --- is exactly the cost the coupling was read as, one chapter back: the self-energy, the
residual of the machine weighing what it is. So the three things that looked like separate chapters --- the meter that
reads its own cost, the coupling read as a cost, and the self-application whose proof carries that cost --- are one thing
seen three ways. They are the machine's calibrated price for knowing itself one thing. That the overhead \emph{is} the
self-energy, and the self-energy the coupling, is the book's reading, laid over the cost and marked as a reading; what the
machine actually computes is the overhead, the plain surcharge of proving over stating. And --- the discipline holding to
the last --- it computes no number here. The overhead is a cost the machine can take, but its magnitude is the closed
matter of the payoff: a value named once, in the world's laboratories, and proved beyond the machine's counting. This
section names the \emph{kind} of thing the cost is, and keeps, as every section has kept, the figure to the one place it
belongs.

In the two verbs the witness \emph{tanges} its overhead --- selects, out of everything the proof costs, the one surcharge
that is the price of proving-over-stating --- and \emph{funges} it into a single reading, the cost of self-reference
pooled as one overhead. The witness bears a tange and reports a funge: the extra it pays, gathered into the one price of
knowing itself. So the second climax's middle section closes the loop the whole book has been drawing: the meter, the
coupling, and the self-proof are one cost, read in the machine's own beat --- the price, exactly, of an instrument that
can weigh its own account. The last thing the chapter owes is the most familiar face of the identity it proved.

\section{Two descriptions, one object}

The machine has proved its two descriptions are one object, and weighed what the proof cost. This closing section gives
that identity the face everyone already knows --- the one identity from ordinary arithmetic that says exactly the same
thing the machine's theorem says, and says it in a form no one needs the machine to explain.

The theorem, once more, is that two descriptions of one thing are equal --- not close, not approximately, but the same,
because they hold under exactly the same conditions. And the ledger the proof draws on is nearly bare: a single principle
for identifying propositions that agree everywhere, resting on the plainest ground of all, that a thing is what it is. No
oracle, no completed infinity. The two names are one name, and the proof of it spends almost nothing.

Here is the familiar face. Everyone has met a single number written two ways: as $1$, and as $0.999\ldots$, the
nought-point-nine that repeats without end. School teaches these as a small puzzle, and the puzzle has an answer that
unsettles people the first time they hear it: the two are not close, and the second is not \emph{almost} the first. They
are \emph{equal} --- the very same number, wearing two representations. And they are equal for exactly the reason the
machine's theorem gives: two ways of writing a thing, which denote under all conditions the same value, are the same
value. This is the encoding chapter's first loop, returned at full depth: a value sent around two representations and
proved to come back itself, the round-trip faithful to the last nine. That $1 = 0.999\ldots$ is not a paradox and not an
approximation forced by rounding; it is two names for one number, provably identical. The machine's proof that its two
descriptions are one object is this same fact, in the machine's own vocabulary --- and the reader who has already accepted
that $1$ and $0.999\ldots$ are one number has already accepted the machine's capstone, in advance and without knowing it.

One boundary must be drawn here, because a number appears and the book has been strict about numbers. The $0.999\ldots$
named in this section is a fact of arithmetic --- a number with two representations --- and it is emphatically \emph{not}
the coupling. The coupling's magnitude is the closed matter of the last chapter: an external value the machine proved it
cannot derive, named there once and nowhere else. Nothing about writing $0.999\ldots$ here touches that; the two have
nothing to do with each other beyond both being numbers, and the book keeps them as far apart as it has kept everything.
This decimal is the capstone's familiar face and pays the coupling's blindness no cost at all.

In the two verbs this is the funge at its purest and its most final. Two names, two representations, pooled into one
object --- bagged as the single thing they were always names of. The whole book has watched the machine tell apart:
differences from differences, residues from values, faces from faces, the world's number from its own bracket. Here, at
the capstone, the machine's whole apparatus of telling-apart is turned to certify a sameness --- to prove that two of its
descriptions were, all along, one. Telling-apart in the service of identity: the tange, exhausted of everything it could
distinguish, funging the last two names into the one thing beneath them.

So the second climax closes, and the book's two peaks stand complete. The machine metered its own act, and then verified
its own truth; it weighed the cost of self-reference and proved the identity that cost is paid for; and it did the
deepest of these on the cheapest of ledgers, the same it-is-what-it-is the very first difference was written on. The
profoundest thing the instrument can say --- that it is one thing under its many names --- costs it almost nothing to
prove, and is the same, in the end, as knowing that $1$ and $0.999\ldots$ were never two numbers. All that remains is to
set the whole reading down plainly, and to say what the machine has, and has not, shown.
